\chapter{Huber's General Information Game}
\label{sec:p15-huber-information-game}

The normal-contamination solution is one instance of a construction
for log-concave prototype densities. Its general explanation uses an
information game on convex classes of measures
\cite[Secs.~5 and 8]{Huber1964RobustLocation}.
Existence of an information minimizer, uniqueness of that minimizer,
and existence of a saddlepoint score are separate conclusions.
Support changes affect the latter two and cannot be discarded by
integrating a logarithmic derivative across disconnected intervals.

\section{The measure space and the variational information}
\label{sec:p15-information-game-foundations}

\begin{hbdefinition}{def:p15-information-game-domain}{Subprobabilities, test scores and two payoffs}
Let \(\mathcal M_{\leq1}\) be the finite nonnegative Borel measures
on \(\mathbb R\) of total mass at most one.
The vague topology is the topology of the maps
\(F\mapsto\int h\,dF\), \(h\in C_c(\mathbb R)\).
Here \(C_c\) denotes continuous compactly supported real functions,
and \(\mathcal D=C_c^1(\mathbb R)\) the continuously differentiable ones.
For \(\psi\in\mathcal D\), define
\[
 A_\psi(F)=\int\psi'\,dF,\quad B_\psi(F)=\int\psi^2\,dF,\quad
 J(\psi,F)=2A_\psi(F)-B_\psi(F).
\]
Set \(K(\psi,F)=A_\psi(F)^2/B_\psi(F)\) when \(B_\psi(F)>0\).
If \(B_\psi=0\), set \(K=\infty\) when \(A_\psi\ne0\), and \(K=0\)
when \(A_\psi=0\). Define
\[
 I(F)=\sup_{\psi\in\mathcal D}K(\psi,F).
\]
In particular \(I(0)=0\).
Nature minimizes \(K\); the statistician maximizes it.
Subprobabilities are part of this abstract game. A measure of mass
less than one is not silently used as the law of an observed real
variable without specifying where its missing mass went.
\end{hbdefinition}

\begin{hbtheorem}{thm:p15-vague-subprobability-compactness}{proved=annex:p15.vague.subprobability.compactness}{The compact domain and the closure of contamination}
The vague topology on \(\mathcal M_{\leq1}\) is metrizable and compact.
Consequently every vaguely closed subset is compact.
For a probability \(G\) and \(0<\varepsilon<1\), the vague closure
of its proper contamination class is exactly
\[
 \{(1-\varepsilon)G+\varepsilon H:H\in\mathcal M_{\leq1}\}
 =\{F\in\mathcal M_{\leq1}:F\geq(1-\varepsilon)G\}.
\]
It is convex and compact. If symmetry is required, the same statement
holds with symmetric measures and symmetric escaping mass.
The proper contamination class itself need not be vaguely closed.
\end{hbtheorem}

\begin{hbtheorem}{thm:p15-information-payoff-bridge}{proved=annex:p15.information.payoff.bridge}{The original ratio, its affine representation and mixtures}
For every \(F\in\mathcal M_{\leq1}\),
\[
 I(F)=\sup_{\psi\in\mathcal D}J(\psi,F).
\]
For nonzero \(F\), this also equals the source's supremum restricted
to \(B_\psi(F)>0\). The conventions at zero denominator and the
zero measure therefore complete its domain without changing that value.
\(I\) is convex and lower semicontinuous for the vague topology.
For \(v_1,v_2>0\) and \(0\leq t\leq1\), the source's Lemma 6 is
\[
 \frac{(tu_1+(1-t)u_2)^2}{tv_1+(1-t)v_2}
 \leq t\frac{u_1^2}{v_1}+(1-t)\frac{u_2^2}{v_2}.
\]
Thus \(K(\psi,\cdot)\) is convex.
If \(C\) is compact and convex and \(\Pi\) is a Borel probability
on \(C\), its barycenter \(\overline F(E)=\int_C F(E)\Pi(dF)\)
belongs to \(C\), and
\[
 K(\psi,\overline F)\leq\int_C K(\psi,F)\Pi(dF).
\]
This is the precise sense in which mixing does not improve Nature's
minimum payoff.
\end{hbtheorem}

\begin{hbtheorem}{thm:p15-finite-information-characterization}{proved=annex:p15.finite.information.characterization}{All of the finite-information regularity conclusion}
For \(F\in\mathcal M_{\leq1}\), \(I(F)<\infty\) exactly when
\(F(dx)=f(x)\,dx\), where \(f\geq0\) is an absolutely continuous
function and
\[
 \int_{\{f>0\}}\frac{f'(x)^2}{f(x)}\,dx<\infty.
\]
Then this integral equals \(I(F)\), \(f'\in L^1(dx)\),
\(f\) is continuous, bounded and of bounded variation, and
\[
 \int|f'|\leq\sqrt{F(\mathbb R)I(F)},\qquad
 \|f\|_\infty\leq\tfrac12\sqrt{F(\mathbb R)I(F)}.
\]
The zero measure is included; \(I(F)=0\) only for that measure.
In particular every nonzero singular measure has infinite information.
The proof supplies the test-function density and bounded-functional
representation needed to establish the result.
\end{hbtheorem}

\begin{hbtheorem}{thm:p15-information-game-value}{proved=annex:p15.information.game.value}{Existence of minimizers and the exact minimax value}
Let \(C\subseteq\mathcal M_{\leq1}\) be nonempty, compact and convex,
and let \(\alpha=\inf_{F\in C}I(F)<\infty\).
An information minimizer exists, and
\[
 \alpha=\inf_{F\in C}\sup_{\psi\in\mathcal D}J(\psi,F)
       =\sup_{\psi\in\mathcal D}\inf_{F\in C}J(\psi,F)
       =\sup_{\psi\in\mathcal D}\inf_{F\in C}K(\psi,F).
\]
There are test scores whose worst payoff approaches \(\alpha\).
An attaining score in \(\mathcal D\) is not asserted.
For example the singleton standard-normal class has value one,
but no compactly supported \(C^1\) score attains it.
\end{hbtheorem}

\section{What support changes do to uniqueness and a saddlepoint}
\label{sec:p15-information-support-correction}

\begin{hbtheorem}{thm:p15-information-equality-components}{proved=annex:p15.information.equality.components}{The exact equality condition replacing blanket uniqueness}
Let \(f,g\) be finite-information densities of subprobabilities, and
\(0<t<1\). The deficit in convexity is
\[
\begin{split}
 &tI(f)+(1-t)I(g)-I(tf+(1-t)g)\\
 &\quad=t(1-t)\int_{\{fg>0\}}
 \frac{fg}{tf+(1-t)g}
       \left(\frac{f'}f-\frac{g'}g\right)^2dx .
\end{split}
\]
Equality holds exactly when the two logarithmic derivatives agree
almost everywhere on their overlap.
Equivalently, on each connected component of \(\{f+g>0\}\),
\(f\) and \(g\) are constant nonnegative multiples of \(f+g\).
Those constants may differ between components and may be zero.

Any two information minimizers in a convex class satisfy this
componentwise condition. If their union of positive supports is
connected, they coincide, provided their common information is positive.
The zero-information minimizer is the zero measure.
In particular a positive-everywhere minimizer is unique.
Convexity and compactness alone do not imply uniqueness.
\end{hbtheorem}

\begin{hbtheorem}{thm:p15-information-first-variation}{proved=annex:p15.information.first.variation}{The support term in the source's variational condition}
Let \(f\) have finite information, put \(S=\{f>0\}\) and
\(\psi=-f'/f\) on \(S\). For a finite-information competitor \(g\), set
\[
 b_g=\int_S\psi^2g,\quad
 r_g=\int_{S^c\cap\{g>0\}}\frac{g'^2}{g}.
\]
If \(b_g=\infty\), the right derivative at zero of
\(I((1-t)f+tg)\) is \(-\infty\).
If \(b_g<\infty\), define \(a_g=-\int_S\psi g'\); this integral
is absolutely convergent, and
that derivative is
\[
                    2a_g-b_g+r_g-I(f).
\]
Thus a finite-information \(f\in C\) minimizes information on a
convex class exactly when \(b_g<\infty\) and
\(2a_g-b_g+r_g\geq I(f)\) for every finite-information \(g\in C\).
Infinite-information competitors cannot lower its value.
The term \(r_g\) must not be omitted when a competitor has new support.
\end{hbtheorem}

\begin{hbdefinition}{def:p15-extended-information-score}{Two precise extensions of the test-score domain}
For a class of finite-information densities, a measurable score
\(\psi\) may be evaluated by
\[
 a_\psi(F)=-\int\psi f',\qquad b_\psi(F)=\int\psi^2f,\qquad
 K_{\rm d}(\psi,F)=a_\psi(F)^2/b_\psi(F),
\]
when these integrals are finite, with the preceding zero-denominator
conventions. This is the source's density-based integration formula.

For a class also containing singular measures, use a real continuous,
locally absolutely continuous \(\psi\), together with a specified
bounded Borel representative \(\chi\) of its weak derivative.
Require \(b_\psi(F)=\int\psi^2dF<\infty\) for every \(F\in C\);
put \(a_\psi(F)=\int\chi\,dF\) and use the same ratio.
The values of \(\chi\) on atoms are part of the strategy.
Ordinary almost-everywhere derivatives of a non-absolutely-continuous
score do not supply this weak-derivative premise.
\end{hbdefinition}

\begin{hbtheorem}{thm:p15-information-saddlepoint-criterion}{proved=annex:p15.information.saddlepoint.criterion}{Exactly what is needed beyond a minimizing density}
Let \(C\) be a convex class and let \(F_*\in C\) minimize information
with \(0<\alpha=I(F_*)<\infty\).
For either admissible score domain above, a pair \((\psi,F_*)\)
is a saddlepoint exactly when
\[
 K(\psi,F_*)=\alpha,\qquad K(\psi,F)\geq\alpha\quad(F\in C).
\]
The first equality requires an efficient score on \(F_*\)'s support;
the second constrains its extension to all other supports.

If every finite-information \(F\in C\) is absolutely continuous
with respect to \(F_*\), the efficient score
\(-f_*'/f_*\) satisfies
\[
             2a_\psi(F)-b_\psi(F)\geq\alpha
\]
on that finite-information class. Its required second moments and
integration identities follow from the first-variation condition;
they are not additional assumed tail bounds.
If \(\psi\) is a continuous real extension of this efficient score
to the whole line, \(f_*\) is positive everywhere and the domination
condition is automatic.
\end{hbtheorem}

\begin{hbtheorem}{thm:p15-information-closure-extension}{proved=annex:p15.information.closure.extension}{The source's extension to singular limits}
Let \(C\) be compact and convex, let its finite-information members
\(C_1\) be vaguely dense, and let \(F_*\) minimize information
with \(0<\alpha<\infty\).
Suppose its efficient score has a real continuous, locally absolutely
continuous extension \(\psi\), whose weak derivative has a representative
\(\chi\) either in \(C_0(\mathbb R)\), or nonnegative, upper
semicontinuous and vanishing at infinity.
Then the candidate has finite \(b_\psi(F)\) for every \(F\in C\) and
\[
                  2a_\psi(F)-b_\psi(F)\geq\alpha .
\]
It is therefore a saddlepoint over the full class in the stated
extended-score domain. The proof establishes every continuity and
semicontinuity passage, including zero-denominator cases.

For proper probability members, the game represents inverse asymptotic
variance only under the location-estimation premises.
In particular, if \(\psi\) is bounded and nondecreasing,
\(\lambda_F(0)=0\), \(\lambda_F'(0)=-d_F<0\) is finite, and
finite-information approximants \(F_n\to F\) satisfy
\(\limsup_n a_\psi(F_n)\leq d_F\), then
\[
 E_F\psi^2/d_F^2\leq1/\alpha .
\]
The corresponding convex M-estimator is asymptotically normal under
the hypotheses of \cref{thm:p15-convex-location-normality}.
\end{hbtheorem}

\section{The general log-concave prototype}
\label{sec:p15-huber-general-prototype}

\begin{hbdefinition}{def:p15-logconcave-prototype}{The source's smooth prototype and its clipped convex potential}
Let \(g\in C^2(\mathbb R)\) be a probability density, with
\(S=\{g>0\}\) an interval, and let \(h=-\log g\) be convex on \(S\).
Extend \(h\) by \(+\infty\) outside \(S\).
For \(k>0\), define
\[
 h_k(x)=\inf_{y\in S}\{h(y)+k|x-y|\},\quad
 q_k(x)=e^{-h_k(x)},\quad Z(k)=\int_{\mathbb R}q_k(x)\,dx.
\]
Let \([t_0,t_1]\), with extended endpoints allowed, be the interval
where \(|h'|\leq k\). Integrals over infinite endpoints have their
usual improper meaning; a missing infinite-endpoint tail contributes zero.
\end{hbdefinition}

\begin{hbtheorem}{thm:p15-huber-general-density}{proved=annex:p15.huber.general.density}{The full tangent-tail construction and normalization}
The interval in the preceding definition is nonempty.
At a finite endpoint, \(h'(t_0)=-k\), \(h'(t_1)=k\), and
\[
 q_k(x)=
 \begin{cases}
 g(t_0)e^{k(x-t_0)},&x\leq t_0,\\
 g(x),&t_0<x<t_1,\\
 g(t_1)e^{-k(x-t_1)},&x\geq t_1 .
 \end{cases}
\]
Omit any branch with an infinite endpoint.
The function \(Z(k)\) is finite, continuous and nonincreasing,
strictly decreasing wherever its value exceeds one, with limits
\(Z(0+)=\infty\), \(Z(\infty)=1\).
For every \(0<\varepsilon<1\), exactly one \(k>0\) solves
\[
 (1-\varepsilon)^{-1}=Z(k)
 =\int_{t_0}^{t_1}g+\frac{g(t_0)+g(t_1)}k,
\]
again dropping infinite-endpoint terms.
Then \(f_*=(1-\varepsilon)q_k\) is a positive \(C^1\) probability
density in the original contamination class, and its location score
\(\psi_k=h_k'\) is the continuous nondecreasing clipping of \(h'\)
at \(-k,k\). It is locally absolutely continuous, with
\(\psi_k'\geq0\) and \(\int\psi_k'\leq2k\).
\end{hbtheorem}

\begin{hbtheorem}{thm:p15-huber-general-minimax}{proved=annex:p15.huber.general.minimax}{The general saddlepoint and its exact statistical interpretation}
For the normalized construction set
\[
 I_*=\int\psi_k^2f_*=(1-\varepsilon)\int\psi_k'g>0.
\]
The excess contaminant is supported on clipped tails, so also
\[
 I_*=(1-\varepsilon)\int\psi_k^2g+\varepsilon k^2.
\]
The density \(f_*\) uniquely minimizes information over both the
proper contamination class and its subprobability vague closure.

For the source's centered class
\[
 \mathcal C_k^0=\{F=(1-\varepsilon)G+\varepsilon H:E_F\psi_k=0\},
\]
take i.i.d.\ observations with law \(F\).
The midpoint convex M-estimator based on \(h_k\) is consistent for zero
and has limiting truncated risk at most \(1/I_*\), even when a
contaminant does not give a population-score derivative.
At \(F_*\) its limit is \(N(0,1/I_*)\).
Consequently its worst limiting truncated risk on \(\mathcal C_k^0\)
is \(1/I_*\), minimal among all translation-equivariant estimators.
The limiting-risk order is exactly that in
\cref{def:p15-huber-symmetric-location-model}.
Whenever \(\lambda_F'(0)=-d_F<0\) is finite, its normal-limit variance
is \(E_F\psi_k^2/d_F^2\leq1/I_*\).
For a symmetric prototype, all symmetric contaminants satisfy the
centering condition.

The case \(\varepsilon=0\) is the nominal model \(G\); finite-information
and normal-limit statements there retain their stated hypotheses.
The positive-contamination construction does not automatically prove
normality for an unbounded nominal score at that endpoint.
\end{hbtheorem}

\begin{hbtheorem}{thm:p15-huber-nonnormal-examples}{proved=annex:p15.huber.nonnormal.examples}{Logistic and asymmetric Gumbel calculations}
For the logistic density \(g(x)=\tfrac14\operatorname{sech}^2(x/2)\),
\(h'(x)=\tanh(x/2)\). For \(0<k<1\),
\[
 t_1=-t_0=\log\frac{1+k}{1-k},\quad
 Z(k)=\frac{1+k^2}{2k},\quad
 \varepsilon=\frac{(1-k)^2}{1+k^2}.
\]
Thus \(k=1/2\) corresponds to \(\varepsilon=1/5\), clipping at
\(\pm\log3\), and
\[
 I_*=\frac{11}{60},\qquad V_*=\frac{60}{11}.
\]
For the sample \((0,0,0,M)\), \(M>\log(21/5)\), its estimate is
\(\log(7/5)\), independent of the outlier magnitude.

For the Gumbel density \(g(x)=\exp(-x-\exp(-x))\), \(h'=1-e^{-x}\).
At \(k=1\), the lower endpoint is \(-\log2\), the upper endpoint
is \(+\infty\), and
\[
 Z(1)=1+e^{-2},\quad
 \varepsilon=\frac1{1+e^2},\quad
 I_*=\frac{1-3e^{-2}}{1+e^{-2}}.
\]
Its score is \(\max(-1,1-e^{-x})\).
The centering condition forces \(H\) to put all its mass on
\((-\infty,-\log2]\). For \((-M,0,0)\), \(M>2\log2\), the estimate
is \(-\log2\). This example exhibits an infinite construction endpoint
and the restriction imposed by centering an asymmetric prototype.
\end{hbtheorem}

\begin{hbtheorem}{thm:p15-information-game-failures}{proved=annex:p15.information.game.failures}{Support, derivative and attainment failures}
For the disjoint densities \(f,g\) of
\cref{thm:p15-fisher-minimizer-nonuniqueness}, every mixture minimizes
information at \(4\pi^2\).
In the density-based score domain, use the efficient score of \(f\)
on its support and the constant one on the support of \(g\).
It has payoff \(4\pi^2\) at \(f\), but zero at \(g\).
Using the corresponding efficient score on both supports instead
gives payoff \(4\pi^2\) at every mixture.
Thus the score extension and Nature's nonuniqueness are both substantive.

For \(\psi(x)={\bf1}_{[0,\infty)}(x)\), its ordinary derivative is
zero almost everywhere, but
\[
 -\int\psi(x)\phi'(x)\,dx=\phi(0)>0
\]
for the normal density \(\phi\). Ordinary a.e. differentiation alone
does not justify the derivative-based payoff.
For \(F=\delta_0\), the variational information is infinite.
\end{hbtheorem}

For the prototype construction, compute \(k\) by bracketing and bisection
of the proved continuous monotone normalizer, then solve the continuous
monotone sample score equation. A bracket for the location is obtained
by expanding until the two score signs differ. The integral and root
equations specify the mathematical computation. For the midpoint
selection, let \(U_n(t)=\sum_i\psi_k(X_i-t)\) and locate both endpoints
\[
 a_n=\sup\{t:U_n(t)>0\},\qquad
 b_n=\inf\{t:U_n(t)<0\}.
\]
Bisection with the corresponding strict sign finds each endpoint;
the estimate is \((a_n+b_n)/2\), also when the zero set is an interval.
Numerical quadrature
requires its own error control.
The general compact-class existence theorem is not an assertion that
an arbitrary infinite-dimensional class has an effective optimization
algorithm. The source's CDF-neighborhood, scale, joint-estimation and
\(M_r\) presentations, including its numerical tables, are supplied
in \cref{sec:p15-huber-remaining-methods}.
The joint clipping-atom limit and the scale comparison permitting
observation signs are supplied in \cref{sec:p15-huber-joint-endpoints}.
